\chapter{Đối Đáp Với Học Trò Lì}
\chapterintro{Không phải lúc nào cũng cần thắng lời}{Cái cần giữ là nhịp lớp và sự tôn trọng, không phải câu cuối cùng.}

\section{Em quên vở rồi ạ}
\begin{manthoai}{Màn thoại}
\textbf{Học sinh:} Em quên vở rồi ạ.\\
\textbf{Thầy/Cô:} Không sao. Hôm nay em nhớ bài bằng đầu nhiều hơn một chút nhé.
\end{manthoai}
\begin{dungkhinào}
Dùng khi không muốn biến chuyện quên đồ thành xung đột đầu tiết.
\end{dungkhinào}
\begin{nhipthem}
Sau câu đáp nhẹ, vẫn cần giao nhiệm vụ thay thế: mượn giấy, ghi nháp, chụp lại bài sau giờ.
\end{nhipthem}
\ngancach

\section{Em chưa làm bài vì tối qua bận}
\begin{manthoai}{Màn thoại}
\textbf{Học sinh:} Em chưa làm bài vì tối qua bận.\\
\textbf{Thầy/Cô:} Thầy cô tin em bận. Giờ mình bận tiếp phần sửa lại cho đỡ thiếu nhé.
\end{manthoai}
\begin{dungkhinào}
Hợp khi muốn giữ mặt mũi cho học sinh nhưng vẫn kéo em đó vào trách nhiệm.
\end{dungkhinào}
\begin{nhipthem}
Có thể thêm: \textit{``Em làm phần ngắn nhất trước, đừng để thiếu cả bài.''}
\end{nhipthem}
\ngancach

\section{Bạn kia cũng chưa làm mà}
\begin{manthoai}{Màn thoại}
\textbf{Học sinh:} Bạn kia cũng chưa làm mà.\\
\textbf{Thầy/Cô:} Thế lát nữa bạn kia sẽ làm việc với thầy cô. Còn bây giờ thầy cô đang nói chuyện với em.
\end{manthoai}
\begin{dungkhinào}
Rất hợp để cắt ngay kiểu kéo người khác vào làm lá chắn.
\end{dungkhinào}
\begin{nhipthem}
Nói bình tĩnh, không dồn ép. Câu này hiệu quả vì nó kéo trọng tâm về đúng người.
\end{nhipthem}
\ngancach

\section{Em chưa hiểu nên chưa làm ạ}
\begin{manthoai}{Màn thoại}
\textbf{Học sinh:} Em chưa hiểu nên chưa làm ạ.\\
\textbf{Thầy/Cô:} Vậy hôm nay mình bắt đầu từ chỗ em chưa hiểu, chứ không bắt đầu từ chỗ em bỏ luôn nhé.
\end{manthoai}
\begin{dungkhinào}
Hợp khi học sinh chống đỡ nhưng trong đó vẫn có phần thật là đang vướng bài.
\end{dungkhinào}
\begin{nhipthem}
Nối ngay bằng một bước rất cụ thể: làm một ý đầu, nêu một chỗ bí, hoặc hỏi bạn cùng bàn trước.
\end{nhipthem}
\ngancach

\section{Em làm rồi nhưng mất tờ giấy đó ạ}
\begin{manthoai}{Màn thoại}
\textbf{Học sinh:} Em làm rồi nhưng mất tờ giấy đó ạ.\\
\textbf{Thầy/Cô:} Không sao. Giờ mình giữ lại điều em còn nhớ trên đầu trước đã.
\end{manthoai}
\begin{dungkhinào}
Dùng khi không muốn sa vào đúng sai của lý do ngay giữa tiết mà vẫn cần kéo học sinh vào trách nhiệm.
\end{dungkhinào}
\begin{nhipthem}
Cho em nêu ý chính hoặc làm một phần ngắn tại chỗ, rồi xử lý chuyện giấy tờ sau giờ nếu cần.
\end{nhipthem}
